% !TeX root = ../初等数学笔记.tex
% 由主文档通过 \input 载入；请编译主文档。
\setcounter{chapter}{22} % 保留原书章号。
\chapter{统计描述与数据解释}
\label{chap:em-statistics}
统计描述把有限数据压缩为可解释的信息，统计推断则进一步使用抽样模型认识总体。压缩必然损失细节，因此应根据问题选择指标，同时检查数据来源、观测单位和抽样方式。本章数值例子均为教学自拟数据。

\section{总体、样本与数据来源}
\emph{总体}（population）是研究涉及的全部对象，\emph{样本}（sample）是实际观测到的部分对象；每个对象上的特征称为变量。抽样前 $X_1,\ldots,X_n$ 是随机变量，观测后得到具体值 $x_1,\ldots,x_n$。不含未知总体参数、仅由样本及已知量计算的量称为\emph{统计量}（statistic）。

普查试图覆盖总体，抽样则只观察部分单位。简单随机抽样、分层抽样等需有明确选择规则；便利样本和自愿回答样本可能系统性偏向某类对象。样本量增大能减少某些随机波动，却不能自动消除选择偏差、漏报或测量偏差。描述“参加调查的人”的结果，不应未经论证扩展为“所有人”的结论。

\section{数据类型、频数与图表}
类别数据表示种类，数值数据表示可进行有意义算术运算的量；给类别编码为 $1,2,3$ 并不自动赋予平均值实际含义。某类出现次数为频数，除以总观测数得到频率，各互斥完备类别的频率和为 $1$。

条形图比较类别，直方图表示分组数值数据。若第 $j$ 组宽度为 $w_j>0$、频率为 $f_j$，应以 $f_j/w_j$ 为高度，使面积代表频率；组宽不同时直接比较柱高会误导。散点图保留两个数值变量的配对关系。分组后一般不能精确恢复原数据，利用组中值计算均值是一种近似，不是原样本均值的恒等变形。

\section{集中趋势与平均数的选择}
样本均值为 $\overline x=n^{-1}\sum_i x_i$。带非负权重 $w_i$ 且总权重为正时，加权均值为 $\sum_iw_ix_i/\sum_iw_i$。把若干组汇总时，组均值应按组样本量加权，而不是无条件取组均值的简单平均。

将数据排序为 $x_{(1)}\le\cdots\le x_{(n)}$，中位数在奇数样本量时取中间一项，在偶数时取中间两项平均；众数为出现次数最多的值，可以有多个。均值利用全部数值，易受极端值影响；中位数主要反映位置顺序。例如自拟数据 $1,2,2,3,12$ 的均值为 $4$，中位数与众数均为 $2$，差异揭示右侧大值的影响，而不说明某一种指标普遍优越。

均值还有最小二乘意义：对任意实数 $a$，
\[
 \sum_i(x_i-a)^2=\sum_i(x_i-\overline x)^2+n(a-\overline x)^2,
\]
交叉项因 $\sum_i(x_i-\overline x)=0$ 消失。因此均值是使平方偏差和最小的唯一常数。绝对偏差和则由中位数最小化：把两端数据成对，$|a-u|+|a-v|$ 在 $a\in[u,v]$ 时最小，所有这些区间的交给出中位位置。

\section{离散程度、标准化与样本方差}
极差为最大值减最小值。对 $n\ge2$，本书的四分位数约定是：先取总体中位数，奇数样本量时去掉该中间项，再分别取上下半数据的中位数作为 $Q_1,Q_3$；四分位距为 $Q_3-Q_1$。其他软件可能采用插值约定，小样本比较时应核对定义。

描述当前数据离散程度可用
\[
 v_n=\frac1n\sum_i(x_i-\overline x)^2,\qquad
 \sqrt{v_n},
\]
分别为以 $n$ 作分母的方差和标准差。$n\ge2$ 时校正样本方差为 $s^2=(n-1)^{-1}\sum_i(x_i-\overline x)^2$。它的统计意义依赖抽样：若 $X_i$ 独立同分布且取值有限、总体均值为 $\mu$、方差为 $\sigma^2$，则
\[
 \sum_i(X_i-\overline X)^2=\sum_i(X_i-\mu)^2-n(\overline X-\mu)^2.
\]
取期望并用 $\operatorname{Var}(\overline X)=\sigma^2/n$，得到期望为 $(n-1)\sigma^2$，故 $\mathbb E[s^2]=\sigma^2$；以 $n$ 为分母则期望为 $(n-1)\sigma^2/n$。此处 $s^2$ 指观测前的随机统计量，其数值在观测后才确定。标准化 $(x_i-\overline x)/s$ 还要求 $s>0$。

\section{相关、线性拟合与推断边界}
对配对数据 $(x_i,y_i)$，记中心化值 $u_i=x_i-\overline x,v_i=y_i-\overline y$。当 $\sum_i u_i^2$ 与 $\sum_i v_i^2$ 均非零时，样本相关系数为
\[
 r=\frac{\sum_i u_iv_i}{\sqrt{(\sum_i u_i^2)(\sum_i v_i^2)}}.
\]
由柯西不等式，$|r|\le1$；等号意味着中心化数据成比例，即全部点落在某条非水平、非竖直直线上。$r=0$ 只表示这种线性指标为零，不保证没有非线性关系；相关也不能单独推出因果。

拟合直线 $y=\alpha+\beta x$ 的最小二乘目标是使 $\sum_i(y_i-\alpha-\beta x_i)^2$ 最小。令 $D=\sum_i u_i^2>0$、$N=\sum_i u_iv_i$，配方得到
\[
 \sum_i(y_i-\alpha-\beta x_i)^2
 =\sum_i v_i^2-\frac{N^2}{D}
  +D\left(\beta-\frac ND\right)^2
  +n(\alpha+\beta\overline x-\overline y)^2.
\]
因此唯一最优解为 $\widehat\beta=N/D,\widehat\alpha=\overline y-\widehat\beta\overline x$。这是一个确定性的拟合结论，不需要先假设误差正态；若要作置信区间或因果解释，则还需额外模型。所有 $x_i$ 相同的 $D=0$ 情形不能唯一确定斜率。

\section{例题、习题与学习检查}
\begin{example}
对自拟数据 $(0,1),(1,2),(2,2)$ 求最小二乘直线。
\end{example}
\begin{solve}
$\overline x=1,\overline y=5/3$，中心化后的 $D=2,N=1$，所以 $\widehat\beta=1/2,\widehat\alpha=7/6$。拟合直线为 $y=7/6+x/2$，三个残差依次为 $-1/6,1/3,-1/6$，其和为零、平方和为 $1/6$。拟合并不保证每个观测点落在直线上。
\end{solve}
\begin{enumerate}
\item 数据 $2,3,3,4,8$ 的均值、中位数与 $v_n$ 是多少？答：$4,3,22/5$。
\item 一组 $10$ 人平均值为 $60$，另一组 $30$ 人平均值为 $80$，合并平均值是多少？答：$75$。
\item 若把所有数据加常数 $c$，均值和方差怎样变化？答：均值加 $c$，方差不变；乘常数 $c$ 时方差乘 $c^2$。
\end{enumerate}

